\section{Experimental setup}
\label{sec:setup}

\subsection{Data}
In our experiments, we used ProcessBench \cite{processbench} as an aggregated dataset containing 3,400 step-by-step solutions to mathematical problems generated by a range of open-weight LLMs, where the first erroneous step of each solution (if any) was manually annotated. The problems themselves span 4 datasets with varying degree of difficulty; these datasets are, from easier to more difficult: GSM8K \cite{gsm8k}, MATH \cite{math}, OlympiadBench \cite{olympiadbench2024} and OmniMATH \cite{Gao2024OmniMATHAU}.
Steps preceding the first error are correct, the annotated step is incorrect and later steps carry no independent label, so we exclude them. To train and evaluate the learned router we split the problems of each dataset into training, validation and test (75/5/20\%), keeping all solutions of the same problem in the same split; all routers are compared on the test split. More details and statistics about the data in Appendix \ref{app:data_stats}.

\subsection{Router Implementations}

We evaluate two implementations of the router $\mathcal{R}$. For the LLM-based router, we consider two models: the open-weight \textit{Qwen3.5-9B}~\cite{qwen3.5} and the proprietary \textit{GPT-5.6 Luna}~\cite{openai2026gpt56}. Both models receive the same zero-shot prompt, provided in Appendix~\ref{app:prompts}, containing the problem, the preceding steps, the target step and the complete artefacts produced by both verifiers for that step (the synthesised Python checker with its execution outcome, and the Lean statement with its proof attempt and compiler outcome). The model returns a single choice among \textit{python}, \textit{lean}, \textit{tie}, \textit{neither} (both verifications are wrong) and \textit{inconclusive} (the step makes no checkable inference); a second call asks the model to audit its own draft and only the reviewed answer is kept. The LLM router thus selects \emph{a posteriori} between the two verifications, whereas the learned router below decides from the step text alone, before any verifier is run.

For the learned router, $\mathcal{R}_{\mathrm{BERT}}$, we use the standard BERT architecture~\cite{devlin-etal-2019-bert} with a binary classification head. The classifier is trained using Binary Cross-Entropy (BCE) to select between the Python and Lean verifiers. Training examples are constructed from ProcessBench~\cite{processbench} by additionally evaluating each step with both verification procedures. For a correct step, we assign the Python label if Python successfully verifies it; otherwise, we assign Lean if Lean succeeds. The same rule is applied to incorrect steps, using the verifier that correctly identifies the step as incorrect. This prioritizes Python when both verifiers are applicable, reducing the use of the more expensive Lean pipeline. A step is routed to Lean when $P(\mathrm{Lean}) \geq 0.5$; implementation and training details, as well as a variant whose threshold is selected on the validation split, are provided in Appendix~\ref{app:bert_details}.

\subsection{Verifier Implementations}
For the Python verifier, we have used Qwen 3 Coder Next \cite{qwen_qwen3_coder_next_tech_report} as an efficient and high-performing option for python code synthesis with only 3B active parameters; the generated checker is executed in a restricted sandbox and the step is accepted only if it returns \texttt{True}.
For the formal language, we have followed previous literature and chose Lean and used the state-of-the-art Goedel-v2-8B-Autoformalizer as $\mathcal{A}$ and Goedel-v2-8B-Prover as $\mathcal{P}$ \cite{goedel2025}; the step is accepted only if a generated proof of the autoformalised statement compiles against Mathlib. Both verifiers are run once on every ProcessBench step; details and outcome statistics are in Appendix~\ref{app:verifiers}.

\subsection{Baselines}

We compare our approach against five baselines reported in ProcessBench~\cite{processbench}. For process reward models, we consider \textit{Skywork-PRM-7B}~\cite{skywork2024} and \textit{Qwen2.5-Math-7B-PRM800K}~\cite{processbench}. We additionally evaluate three LLM-based critics: \textit{QwQ-32B-Preview}~\cite{qwq2024}, \textit{GPT-4o-0806}~\cite{openai2024gpt4o}, and \textit{o1-mini}~\cite{openai2024o1mini}. These models are evaluated as step-level critics following the ProcessBench protocol, providing a comparison with existing approaches for identifying errors in mathematical reasoning.

\subsection{Evaluation}
We evaluate at the step level: a step is predicted correct when the verifier selected by the router returns \texttt{True} and incorrect otherwise, so that autoformalisation and prover failures count as rejections. Since about 90\% of the annotated steps are correct, we report \emph{balanced accuracy}, the mean of the recall on correct and on incorrect steps, together with plain accuracy. Routers are compared on the test split of the learned router, restricted to the steps for which every method produces a verdict (Appendix~\ref{app:evaluation}). The ProcessBench baselines are instead reported with the solution-level F1 of the original paper, computed on the whole benchmark, and are therefore only indicative.
